\documentclass[11pt]{article}

\usepackage[a4paper, margin=2.5cm]{geometry}
\usepackage{amsmath, amssymb}
\usepackage{exsheets}

\title{Taller II}
\author{Curso: Matemática Estructural, Carla Gonzalez Mina}
\date{Due: Lunes 31 Agosto, 2026}

\SetupExSheets{
  question/name = Ejercicio,
  solution/name = Solución
}

\newcommand{\N}{\mathbb{N}}
\newcommand{\Pow}{\mathcal{P}}

\begin{document}
\maketitle

\begin{question}
Sean $A,B,C$ y $D$ conjuntos. Compare las siguientes parejas de conjuntos. Si son iguales, pruebe su afirmación por doble inclusión o por álgebra de conjuntos. Si uno de los dos conjuntos está contenido propiamente en el otro, pruebe la inclusión y muestre con un ejemplo que dicha inclusión es propia. Si son incomparables, muéstrelo con un contraejemplo.

\begin{enumerate}
  \item[(a)] $\Pow(A)\cap\Pow(B)$ y $\Pow(A\cap B)$.
  \item[(b)] $(A\triangle B)\times C$ y $(A\times C)\triangle(B\times C)$.
  \item[(c)] $(A\triangle B)\cup(C\triangle D)$ y $(A\cup C)\triangle(B\cup D)$.
  \item[(d)] $A\setminus(B\cup C)$ y $(A\setminus B)\setminus C$.
\end{enumerate}
\end{question}

\begin{solution}
\textbf{(a)} Voy a demostrar por doble inclusión que
\[
\boxed{\Pow(A)\cap\Pow(B)=\Pow(A\cap B)}.
\]
Recuerdo que, para cualesquiera conjuntos $X$ y $Y$,
\[
X\in\Pow(Y)\Longleftrightarrow X\subseteq Y.
\]

$(\subseteq)$ Sea $X\in\Pow(A)\cap\Pow(B)$. Por definición de intersección,
\[
X\in\Pow(A)\qquad\text{y}\qquad X\in\Pow(B).
\]
Por definición del conjunto potencia, $X\subseteq A$ y $X\subseteq B$.
Esto quiere decir que, para todo elemento $x$,
\[
x\in X\Longrightarrow x\in A\text{ y }x\in B.
\]
Por definición de intersección, $x\in A\cap B$. Por consiguiente,
$X\subseteq A\cap B$, y nuevamente por definición del conjunto potencia,
\[
X\in\Pow(A\cap B).
\]
Como $X$ fue arbitrario, he demostrado que
\[
\Pow(A)\cap\Pow(B)\subseteq\Pow(A\cap B).
\]

$(\supseteq)$ Sea ahora $X\in\Pow(A\cap B)$. Entonces
$X\subseteq A\cap B$. Como $A\cap B\subseteq A$ y $A\cap B\subseteq B$,
por transitividad de la inclusión obtengo
\[
X\subseteq A\qquad\text{y}\qquad X\subseteq B.
\]
Por definición del conjunto potencia, $X\in\Pow(A)$ y $X\in\Pow(B)$; por
definición de intersección,
\[
X\in\Pow(A)\cap\Pow(B).
\]
Esto demuestra que $\Pow(A\cap B)\subseteq\Pow(A)\cap\Pow(B)$.
Concluyo por doble inclusión que
\[
\Pow(A)\cap\Pow(B)=\Pow(A\cap B).
\]

\textbf{(b)} También se cumple la igualdad
\[
\boxed{(A\triangle B)\times C=(A\times C)\triangle(B\times C)}.
\]
La demostraré siguiendo una pareja ordenada arbitraria. Sea $(x,y)$ una
pareja cualquiera. Entonces
\begin{align*}
(x,y)\in(A\triangle B)\times C
&\Longleftrightarrow x\in A\triangle B\text{ y }y\in C\\
&\Longleftrightarrow
  [(x\in A\text{ y }x\notin B)\text{ o }(x\in B\text{ y }x\notin A)]
  \text{ y }y\in C.
\end{align*}
Distribuyendo la condición $y\in C$ entre los dos casos, esto equivale a
\begin{align*}
&[(x\in A,\ y\in C)\text{ y no ocurre }(x\in B,\ y\in C)]\\
&\qquad\text{o}\quad
[(x\in B,\ y\in C)\text{ y no ocurre }(x\in A,\ y\in C)].
\end{align*}
Por definición del producto cartesiano, la condición anterior dice que
$(x,y)$ pertenece exactamente a uno de $A\times C$ y $B\times C$. Por
definición de diferencia simétrica,
\[
(x,y)\in(A\times C)\triangle(B\times C).
\]
Todos los pasos anteriores son equivalencias, de modo que también se obtiene
la implicación recíproca. Por tanto, para toda pareja $(x,y)$,
\[
(x,y)\in(A\triangle B)\times C
\Longleftrightarrow
(x,y)\in(A\times C)\triangle(B\times C).
\]
Por extensionalidad concluyo la igualdad requerida.

\textbf{(c)} Aquí, en general, lo que tengo es
\[
\boxed{(A\cup C)\triangle(B\cup D)
\subseteq(A\triangle B)\cup(C\triangle D)},
\]
y noto que la inclusión puede ser propia.

Sea $x\in(A\cup C)\triangle(B\cup D)$. Por definición de diferencia
simétrica, tengo exactamente dos casos.

\emph{Caso 1.} Supongamos que
\[
x\in(A\cup C)\setminus(B\cup D).
\]
Entonces $x\in A\cup C$ y $x\notin B\cup D$. La segunda condición implica
$x\notin B$ y $x\notin D$, mientras que la primera implica $x\in A$ o
$x\in C$. Si $x\in A$, como $x\notin B$, se tiene
$x\in A\setminus B\subseteq A\triangle B$. Si $x\in C$, como $x\notin D$,
se tiene $x\in C\setminus D\subseteq C\triangle D$. En cualquiera de las dos
posibilidades,
\[
x\in(A\triangle B)\cup(C\triangle D).
\]

\emph{Caso 2.} Supongamos que
\[
x\in(B\cup D)\setminus(A\cup C).
\]
Entonces $x\in B\cup D$ y $x\notin A\cup C$. Por las leyes de De Morgan,
$x\notin A$ y $x\notin C$. Además, $x\in B$ o $x\in D$. Si $x\in B$,
entonces $x\in B\setminus A\subseteq A\triangle B$; si $x\in D$, entonces
$x\in D\setminus C\subseteq C\triangle D$. De nuevo,
\[
x\in(A\triangle B)\cup(C\triangle D).
\]

Como los dos casos posibles llevan al mismo conjunto, queda demostrada la
inclusión
\[
(A\cup C)\triangle(B\cup D)
\subseteq(A\triangle B)\cup(C\triangle D).
\]

Para ver que puede ser propia, tomo
\[
A=D=\{1\},\qquad B=C=\varnothing.
\]
Entonces
\[
(A\triangle B)\cup(C\triangle D)=\{1\},
\]
mientras que
\[
(A\cup C)\triangle(B\cup D)=\{1\}\triangle\{1\}=\varnothing.
\]
En particular, el elemento $1$ pertenece al conjunto de la derecha y no al
conjunto de la izquierda. Por tanto, la inclusión anterior puede ser propia.

\textbf{(d)} Aquí también se cumple la igualdad
\[
\boxed{A\setminus(B\cup C)=(A\setminus B)\setminus C}.
\]
Pensando en $A,B,C$ como subconjuntos de un universo, aplico las leyes de De Morgan:
\begin{align*}
A\setminus(B\cup C)
&=A\cap(B\cup C)^c
&&\text{(definición de diferencia)}\\
&=A\cap(B^c\cap C^c)
&&\text{(ley de De Morgan)}\\
&=(A\cap B^c)\cap C^c
&&\text{(asociatividad de la intersección)}\\
&=(A\setminus B)\cap C^c
&&\text{(definición de diferencia)}\\
&=(A\setminus B)\setminus C
&&\text{(definición de diferencia).}
\end{align*}
Cada paso es una igualdad válida para conjuntos, de modo que la primera y la
última expresión representan el mismo conjunto.
\end{solution}

\begin{question}
Decidan si la diferencia simétrica es asociativa o no:
\[
A\triangle(B\triangle C)=(A\triangle B)\triangle C.
\]
Si es asociativa, demuéstrelo. Si no, dé un contraejemplo.

¿Se les ocurre una manera fácil de describir $A\triangle(B\triangle C)$? ¿Funciona para $(A\triangle D)\triangle(B\triangle C)$?
\end{question}

\begin{solution}
La diferencia simétrica sí es asociativa:
\[
\boxed{A\triangle(B\triangle C)=(A\triangle B)\triangle C}.
\]
Por definición, $x\in X\triangle Y$ si y sólo si $x$ pertenece exactamente a
uno de $X$ y $Y$. Defino
\[
E:=\{x:x\text{ pertenece a un número impar de los conjuntos }A,B,C\}.
\]
Para tres conjuntos, esto significa que $x$ pertenece exactamente a uno de ellos o a los tres. Escrito de forma explícita,
\begin{align*}
E={}&[A\setminus(B\cup C)]\cup[B\setminus(A\cup C)]
\cup[C\setminus(A\cup B)]\\
&\cup(A\cap B\cap C).
\end{align*}

Voy a demostrar que ambos lados son iguales a $E$.

Primero, sea $x\in A\triangle(B\triangle C)$. Por definición de diferencia
simétrica ocurre exactamente uno de los siguientes casos:
\begin{enumerate}
  \item $x\in A$ y $x\notin B\triangle C$;
  \item $x\notin A$ y $x\in B\triangle C$.
\end{enumerate}
En el primer caso, $x$ pertenece a $A$ y pertenece a cero o a dos de los
conjuntos $B,C$; en total pertenece a uno o a tres de $A,B,C$. En el segundo
caso, no pertenece a $A$ y pertenece exactamente a uno de $B,C$; en total
pertenece a uno de los tres conjuntos. En cualquier caso, $x\in E$. Esto prueba
que
\[
A\triangle(B\triangle C)\subseteq E.
\]

Para la inclusión contraria, sea $x\in E$. Si $x\in A$, entonces el número de
pertenencias de $x$ entre $B,C$ debe ser par: pertenece a ambos o a ninguno.
Por tanto, $x\notin B\triangle C$, y así
$x\in A\triangle(B\triangle C)$. Si $x\notin A$, para que el número total de
pertenencias sea impar, $x$ debe pertenecer exactamente a uno de $B,C$; esto es,
$x\in B\triangle C$. De nuevo, $x\in A\triangle(B\triangle C)$. Así,
\[
E\subseteq A\triangle(B\triangle C),
\]
y por doble inclusión,
\[
A\triangle(B\triangle C)=E.
\]

Ahora considero el otro lado. Un elemento $x$ pertenece a
$(A\triangle B)\triangle C$ si y sólo si ocurre uno de estos dos casos:
\begin{enumerate}
  \item $x\in A\triangle B$ y $x\notin C$;
  \item $x\notin A\triangle B$ y $x\in C$.
\end{enumerate}
En el primer caso, $x$ pertenece exactamente a uno de $A,B$ y no pertenece a
$C$, de modo que tiene una pertenencia en total. En el segundo caso, pertenece
a cero o a dos de $A,B$ y también pertenece a $C$, de modo que tiene una o tres
pertenencias en total. Recíprocamente, si $x$ pertenece a un número impar de
$A,B,C$, al separar la pertenencia a $C$ se obtiene necesariamente uno de esos
dos casos. Por consiguiente,
\[
(A\triangle B)\triangle C=E.
\]
Como ambos lados son iguales a $E$, por transitividad de la igualdad concluyo
que
\[
A\triangle(B\triangle C)=(A\triangle B)\triangle C.
\]

La misma descripción funciona con cuatro conjuntos:
\[
\boxed{(A\triangle D)\triangle(B\triangle C)
=\{x:x\text{ pertenece a un número impar de }A,B,C,D\}}.
\]
Aquí, ``un número impar'' significa exactamente uno o exactamente tres.
En efecto, cada diferencia simétrica registra la paridad de las pertenencias:
dos apariciones se cancelan y una aparición permanece. Por eso agrupar los
cuatro conjuntos en parejas no cambia la paridad total.
\end{solution}

\begin{question}
Para este punto, recuerde cómo construíamos la biyección entre conjuntos dadas dos inyecciones $f$ y $g$: mirábamos la cadena más larga de ``ancestros'' y si era par, usábamos $f$ y si era impar usábamos $g^{-1}$. (Si no la recuerda y no entiende lo que dije, pregúntele a sus compañeros, o a Antonio, o a mí).

\begin{enumerate}
\item[(a)] Sea $A:=\{a_1,a_2,\ldots,a_n,\ldots\}$ y $B:=\{b_1,b_2,\ldots,b_n,\ldots\}$. Sea $f:A\to B$ definida por $f(a_i)=b_{i+1}$ y $g:B\to A$ definida por $f(b_i)=a_{i+1}$. Para cada $n\in\N$, decida si la ``cadena de ancestros'' de $a_n$ es par o impar, y si para la biyección $\phi$ que construimos tenemos que $\phi(a_n)=f(a_n)$ o si $\phi(a_n)=g^{-1}(a_n)$.

\item[(b)] Sea $f:\N\times\N\to\N$ la función $f((m,n))=2^m3^n$ y sea $g:\N\to\N\times\N$ la función $g(n)=(n,0)$. De acuerdo al Teorema de Schröder--Bernstein, existe un $\phi:\N\times\N\to\N$ que es isomorfismo.

\begin{itemize}
  \item Recuerde la definición de $\phi$ (puede ser la que vimos o la del libro, pero debe ser explícita la que usa). Puede usar en su definición términos como ``$h(a)$'', ``ancestro de $a$'', o, si está siguiendo el texto de Forero, ``profundidad'' de $a$.
  \item De acuerdo a la definición anterior, halle $\phi((256,0))$.
  \item De acuerdo a la definición anterior, halle $\phi((5,3))$.
  \item De acuerdo a la definición anterior, halle $\phi((n,m))$ para $m\neq1$.
  \item De acuerdo a la definición anterior, halle $\phi((5,0))$.
  \item De acuerdo a la definición anterior, halle
  \[
  \phi\!\left(\left(2^{2^{3^2}},0\right)\right).
  \]
\end{itemize}
\end{enumerate}
\end{question}

\begin{solution}

\textbf{Observación sobre el enunciado.} En el inciso (a), la expresión
``$g:B\to A$ definida por $f(b_i)=a_{i+1}$'' debe leerse como
\[
g(b_i)=a_{i+1},
\]
porque $b_i\in B$ y la función que tiene dominio $B$ es $g$. En el inciso (b),
la condición impresa $m\neq1$ parece ser una errata; la condición compatible
con $g(n)=(n,0)$ es $m\neq0$. Más adelante explico qué sucede si se conserva
la lectura literal.

\textbf{Definición de la biyección.} Sean $f:A\to B$ y $g:B\to A$ inyectivas.
Dado $y\in A\cup B$, digo que $x$ es el ancestro de $y$ si
\[
y\in B\text{ y }f(x)=y,
\qquad\text{o si}\qquad
y\in A\text{ y }g(x)=y.
\]
Como $f$ y $g$ son inyectivas, cada elemento tiene a lo sumo un ancestro inmediato.
Digo que $a\in A$ tiene profundidad $k$ si su cadena
\[
a=x_0,x_1,\ldots,x_k
\]
tiene $k$ ancestros propios y $x_k$ no tiene ancestro. Si la cadena no termina,
la profundidad es infinita. Con esto defino
\begin{align*}
X_{\mathrm{par}}&:=\{a\in A:a\text{ tiene profundidad par}\},\\
X_{\mathrm{impar}}&:=\{a\in A:a\text{ tiene profundidad impar}\},\\
X_{\infty}&:=\{a\in A:a\text{ tiene profundidad infinita}\}.
\end{align*}
Con esto, la biyección queda
\[
\boxed{
\phi(a)=
\begin{cases}
f(a),&a\in X_{\mathrm{par}}\cup X_{\infty},\\
g^{-1}(a),&a\in X_{\mathrm{impar}}.
\end{cases}}
\]
La segunda rama está bien definida: si $a$ tiene profundidad impar, su cadena
tiene por lo menos un ancestro; como $a\in A$, ese ancestro es un elemento
$b\in B$ tal que $g(b)=a$. Por tanto, existe $g^{-1}(a)=b$.

\textbf{(a)} Para $n\geq1$, el primer ancestro de $a_n$ es $b_{n-1}$ cuando
$n\geq2$, porque
\[
g(b_{n-1})=a_n.
\]
El ancestro de $b_{n-1}$ es $a_{n-2}$ cuando $n\geq3$, porque
\[
f(a_{n-2})=b_{n-1}.
\]
Continuando alternadamente con $g^{-1}$ y $f^{-1}$, obtengo la cadena
\[
a_n,\ b_{n-1},\ a_{n-2},\ b_{n-3},\ldots
\]
Si $n$ es impar, la cadena termina en $a_1$; este elemento no tiene ancestro
porque no existe un índice $0$ con $g(b_0)=a_1$. Si $n$ es par, la cadena
termina en $b_1$, que tampoco tiene ancestro porque no existe $a_0$ con
$f(a_0)=b_1$. En ambos casos hay exactamente $n-1$ ancestros propios, de modo
que la profundidad de $a_n$ es $n-1$.

Si $n$ es impar, $n-1$ es par y uso $f$:
\[
\phi(a_n)=f(a_n)=b_{n+1}.
\]
Si $n$ es par, $n-1$ es impar y uso $g^{-1}$. Como
$g(b_{n-1})=a_n$, se tiene
\[
\phi(a_n)=g^{-1}(a_n)=b_{n-1}.
\]
Por tanto,
\[
\boxed{
\phi(a_n)=
\begin{cases}
b_{n+1},&n\text{ impar},\\
b_{n-1},&n\text{ par}.
\end{cases}}
\]

\textbf{(b)} Aquí $0\in\N$. Antes de aplicar la construcción, verifico las
inyecciones. Sean $(m,n),(r,s)\in\N\times\N$ y supongamos que
\[
f((m,n))=f((r,s)).
\]
Entonces $2^m3^n=2^r3^s$. Por la unicidad de la factorización en números
primos, los exponentes de $2$ y de $3$ deben coincidir: $m=r$ y $n=s$. Así,
$(m,n)=(r,s)$, lo cual demuestra que $f$ es inyectiva. Para $g$, si
$g(k)=g(\ell)$, entonces
\[
g(k)=g(\ell)\Longrightarrow(k,0)=(\ell,0)\Longrightarrow k=\ell.
\]
Por tanto, $g$ también es inyectiva.

Para $(256,0)$, construyo la cadena
\[
(256,0),\ 256,\ (8,0),\ 8,\ (3,0),\ 3,\ (0,1).
\]
Cada paso se justifica mediante
\[
g(256)=(256,0),\quad f((8,0))=2^8=256,\quad g(8)=(8,0),
\]
\[
f((3,0))=2^3=8,\quad g(3)=(3,0),\quad f((0,1))=3.
\]
La cadena termina en $(0,1)$ porque este elemento no está en la imagen de $g$:
todo elemento de $g(\N)$ tiene segunda coordenada $0$. Hay seis ancestros
propios, de modo que la profundidad es $6$, que es par. Por definición de
$\phi$, uso $f$:
\[
\boxed{\phi((256,0))=f((256,0))=2^{256}}.
\]

Como $(5,3)$ tiene segunda coordenada $3\neq0$, no pertenece a la imagen de
$g$, que es $\N\times\{0\}$. Por tanto, no tiene ancestros y su profundidad es
$0$, que es par. En consecuencia,
\[
\boxed{\phi((5,3))=f((5,3))=2^5 3^3=864}.
\]

En general, si $m\neq0$, la pareja $(n,m)$ no está en la imagen de $g$, pues
ningún valor de $g$ tiene segunda coordenada distinta de $0$. Así, $(n,m)$ no
tiene ancestros, su profundidad es $0$ y corresponde usar $f$:
\[
\boxed{\phi((n,m))=2^n3^m\qquad(m\neq0)}.
\]
Si se conserva literalmente la condición $m\neq1$, esta fórmula resuelve
los casos $m\geq2$, pero el caso $m=0$ también satisface $m\neq1$ y requiere
examinar su cadena por separado.

Para $(5,0)$, el primer ancestro es $5$, porque $g(5)=(5,0)$. La cadena es
\[
(5,0),\ 5.
\]
Para continuar sería necesario encontrar $(r,s)\in\N\times\N$ tal que
$f((r,s))=5$, es decir, $2^r3^s=5$. Esto es imposible por la factorización
prima de $5$. Por tanto, $5$ no tiene ancestro y la profundidad de $(5,0)$ es
$1$, que es impar. Debo usar $g^{-1}$:
\[
\boxed{\phi((5,0))=g^{-1}((5,0))=5}.
\]

Por último, como $3^2=9$ y $2^9=512$, la cadena queda
\[
(2^{512},0),\ 2^{512},\ (512,0),\ 512,\ (9,0),\ 9,\ (0,2).
\]
Los pasos se justifican porque
\[
g(2^{512})=(2^{512},0),\quad f((512,0))=2^{512},
\]
\[
g(512)=(512,0),\quad f((9,0))=2^9=512,
\]
\[
g(9)=(9,0),\quad f((0,2))=3^2=9.
\]
La cadena termina en $(0,2)$ porque su segunda coordenada no es $0$. Tiene seis
ancestros propios; por tanto, su profundidad es par y corresponde usar $f$:
\[
\boxed{
\phi\!\left(\left(2^{2^{3^2}},0\right)\right)
=2^{2^{2^{3^2}}}=2^{2^{512}}}.
\]
\end{solution}

\begin{question}
Sea $f:A\to B$ una función. Para $X\subset A$ defina
\[
f(X):=\{y\in B:\exists x,\ x\in X\text{ y }f(x)=y\}.
\]
De igual manera, para $Y\subset B$ defina
\[
f^{-1}(Y):=\{x\in A:\exists y,\ y\in Y\text{ y }f(x)=y\}.
\]

\begin{enumerate}
  \item[(a)] Sea $A=B=\N$, y sea $f(x)=x^2$. Halle $f(A)$.
  \item[(b)] Halle un contraejemplo para la siguiente afirmación. Si $f$ es una función de $A$ en $B$ y $X\subset A$, entonces $f^{-1}(f(X))=X$.
  \item[(c)] Demuestre, en cambio, la siguiente afirmación. Si $f$ es una función de $A$ en $B$ y $Y\subset B$, entonces $f(f^{-1}(Y))=Y$.
\end{enumerate}
\end{question}

\begin{solution}
\textbf{Observación sobre el enunciado.} Antes de empezar, noto que el inciso (c) no es
verdadero para una función arbitraria. La identidad que sí vale sin hipótesis
adicionales es
\[
f(f^{-1}(Y))=Y\cap f(A).
\]

\textbf{(a)} Por definición de imagen,
\[
f(A)=\{f(n):n\in\N\}=\{n^2:n\in\N\}.
\]
Como en este taller uso $0\in\N$, obtengo
\[
\boxed{f(A)=\{0,1,4,9,16,25,36,\ldots\}}.
\]

\textbf{(b)} Busco un contraejemplo. Tomo
\[
A=\{a,b,c\},\qquad B=\{0,1\},
\]
con $f(a)=0$, $f(b)=0$ y $f(c)=1$. Esta regla define una función
$f:A\to B$, pues a cada elemento de $A$ le asigna exactamente un elemento de
$B$. Elijo $X=\{b\}$. Como el único elemento de $X$ es $b$ y $f(b)=0$,
\[
f(X)=\{f(b)\}=\{0\}.
\]
Ahora,
\[
f^{-1}(f(X))=f^{-1}(\{0\})
=\{x\in A:f(x)=0\}=\{a,b\}.
\]
Como $\{a,b\}\neq\{b\}=X$, se obtiene
\[
f^{-1}(f(X))\neq X.
\]
Este ejemplo refuta la afirmación universal.

La inclusión que siempre se cumple es $X\subseteq f^{-1}(f(X))$: si $x\in X$,
entonces $f(x)\in f(X)$ por definición de imagen, y por definición de imagen
inversa esto implica que $x\in f^{-1}(f(X))$. La igualdad puede fallar porque
un elemento exterior a $X$ puede tener la misma imagen que un elemento de $X$,
como ocurre con $a$ y $b$ en el contraejemplo.

\textbf{(c)} Tal como está escrito, el enunciado no se puede demostrar. Sea
\[
A=\{a\},\qquad B=\{0,1\},\qquad f(a)=0,
\]
y sea $Y=\{1\}$. Entonces
\[
f^{-1}(Y)=\varnothing,
\qquad
f(f^{-1}(Y))=\varnothing\neq\{1\}=Y.
\]

Ahora demuestro la identidad correcta. Sea $y\in B$. Entonces
\begin{align*}
y\in f(f^{-1}(Y))
&\Longleftrightarrow
\text{existe }x\in f^{-1}(Y)\text{ tal que }f(x)=y\\
&\Longleftrightarrow
\text{existe }x\in A\text{ tal que }f(x)=y\text{ y }f(x)\in Y\\
&\Longleftrightarrow y\in f(A)\text{ y }y\in Y\\
&\Longleftrightarrow y\in Y\cap f(A).
\end{align*}
Como esto vale para todo $y$, por extensionalidad concluyo que
\[
\boxed{f(f^{-1}(Y))=Y\cap f(A)}.
\]
De aquí se sigue que $f(f^{-1}(Y))=Y$ cuando $Y\subseteq f(A)$; en particular, la
igualdad pedida se cumple para todo $Y\subseteq B$ si $f$ es sobreyectiva.
En efecto, si $Y\subseteq f(A)$, entonces $Y\cap f(A)=Y$. Si $f$ es
sobreyectiva, $f(A)=B$, y como $Y\subseteq B$, automáticamente
$Y\subseteq f(A)$.
\end{solution}

\clearpage
\section*{Soluciones}
\printsolutions

\end{document}
